\par\clearpage
\noindent\makebox[0pt][l]{\fontsize{13}{16}\selectfont 416}\hfill{\small\textbf{NACHLASS.}}\hfill\null\par\medskip
\noindent dings wichtig, die großen südlichen Dreiecke noch in mehrere kleinere
zu unterteilen und durch Messungen an neu eingefügten Standpunkten sichere
und ausreichende Grundlagen für jene Aufnahmen zu gewinnen. Die Kosten \textls[120]{dieser}
Arbeiten habe ich bei der obigen Schätzung nicht berücksichtigen können:
Ihre Veranschlagung wäre fast noch schwieriger, aber auf jeden Fall können
sie im Vergleich zu den Kosten neuer großer Triangulationen in den noch
nicht erfassten Landesteilen nur gering sein.

Dass es übrigens in Zukunft wünschenswert sein wird, die Ergebnisse zur Lage
aller genau bestimmten Punkte, sobald sie ein geschlossenes Ganzes bilden,
\textls[120]{öffentlich} bekannt zu machen, brauche ich nicht zu bemerken.
Für den eigentlich rein wissenschaftlichen Teil der bisherigen Messungen
versteht sich dies ohnehin von selbst.\quad\mbox{.\ .\ .\ .\ .\ .}

\bigskip
\begin{center}
[Aus dem Bericht vom 26.~Juni 1828, »die Fortsetzung der Gradmessungs-\par
arbeiten betreffend.«]
\end{center}

\par\noindent\hspace*{\parindent}\mbox{.\ .\ .\ .\ .\ .}\quad Da die Detailaufnahme der noch nicht vermessenen Landesteile
auf den trigonometrischen Arbeiten aufbauen soll und beide Aufgaben
hinsichtlich der einzusetzenden Geldmittel voneinander abhängen werden,
war zunächst ein ungefährer Überschlag der Gesamtkosten erforderlich. Nach
einer mir vom Geheimen Kabinettsrat \textsc{Hoppenstedt} mitgeteilten Notiz
würde der Flächeninhalt der im Detail aufzunehmenden Landesteile etwa 144
Quadratmeilen betragen. Die Kosten der Detailaufnahme durch Generalstabsoffiziere
werden auf 200---250 Taler für jede Quadratmeile geschätzt, wozu noch etwa
21 Taler für die Anfertigung von 4 Exemplaren der Karte hinzuzurechnen wären.
Würden also insgesamt 250 Taler je Quadratmeile angesetzt, so würden diese
Kosten etwa 36000 Taler und folglich einschließlich der Triangulationskosten
in den von der Gradmessung noch nicht erfassten Landesteilen gegen 50000 Taler
betragen. Hinzuzurechnen wären wohl noch ein oder ein paar Tausend Taler für
die Arbeiten, die nötig sein werden, um innerhalb der großen südlichen
Dreiecke der Gradmessung eine hinreichend große Anzahl
\par
\clearpage
\noindent\null\hfill{\small\textbf{ZUR HANNOVERSCHEN TRIANGULATION.}}\hfill\makebox[0pt][r]{\fontsize{13}{16}\selectfont 417}\par\medskip
\noindent fester Punkte für die Detailaufnahme festzulegen, worüber ich mich
bereits früher in der im November vorigen Jahres eingereichten Eingabe
ausführlicher geäußert habe. Wenn daher für diese Aufgaben jährlich
tatsächlich 5000 Taler eingesetzt werden können (was außer von den Geldmitteln
auch von der ständigen Verfügbarkeit des Personals abhängen wird), wäre bis
zum vollständigen Abschluss ungefähr ein Zeitraum von 10 Jahren erforderlich.

Nach einem von mir selbst früher angestellten, wenn auch vielleicht weniger
zuverlässigen Überschlag wäre der Flächeninhalt der im Detail aufzunehmenden
Landesteile 173 Quadratmeilen. Danach würde der veranschlagte Kostenaufwand
etwa 7000 Taler höher ausfallen und die Arbeit 1 bis 2 Jahre länger dauern.
Ebenso wäre, wenn die festgesetzten 5000 Taler nur als jeweilige Obergrenze
und nicht als Betrag zu betrachten sind, der tatsächlich jedes Jahr oder im
Durchschnitt einzusetzen ist, eine entsprechende Verlängerung der Dauer des
gesamten Vorhabens die Folge.

Was die Verteilung der trigonometrischen Arbeiten auf die einzelnen Jahre
betrifft, so dürfte es, um diese Arbeiten als Grundlage der Detailaufnahme
schneller abschließen zu können, ratsam sein, anfangs den größeren Teil der
Geldmittel für sie und den kleineren für die Detailaufnahme einzusetzen,
vielleicht im Verhältnis von \GaussIXMathCopy{2/3}{\(\frac{2}{3}\)} zu \GaussIXMathCopy{1/3}{\(\frac{1}{3}\)}. Für letztere
wäre dann in den späteren Jahren, wenn überall eine sichere Grundlage
vorhanden ist, wenn die Arbeiter nach und nach immer geübter sind und wenn
die Geldmittel allein dafür eingesetzt werden können, ein umso rascherer
Fortschritt gewiss. Was die Reihenfolge der trigonometrischen Arbeiten angeht,
besteht beim gegenwärtigen Stand der Dinge kein weiterer Grund, eine Arbeit
der anderen vorzuziehen, als dass darauf geachtet werden muss, die
gleichzeitige Detailaufnahme stets mit genügend zuverlässig bestimmten
Punkten zu versorgen, damit sie niemals aus Mangel an solchen Punkten ins
Stocken zu geraten droht. Eine genauere Festlegung dürfte
\textls[120]{für den Augenblick} teils undurchführbar, teils unnötig, teils nicht
einmal ratsam sein, weil ein gewisser Grad an Freiheit, das den jeweiligen
Umständen nach Zweckmäßigste zu bearbeiten, einem rascheren und besseren
Fortschritt nur förderlich sein kann.
\par\noindent\mbox{.\ .\ .\ .\ .\ .}\par
\medskip\noindent\hspace*{0.03\linewidth}{\footnotesize IX.}\hfill\makebox[0.18\linewidth][l]{53}\par
